\section{Optimization on common log canonical slices}\label{sec:E}

Throughout this section, $x\in S$ is a fixed normal complex affine germ of
positive dimension $n$, $S$ is $\Q$-Gorenstein and klt, and
$R=\mathcal O_{S,x}$, $\mathfrak m=\mathfrak m_x$, and $A=A_S$.
A valuation \emph{over $x$} is nonnegative on $R$; its centre may be a
proper generization of $x$. A valuation \emph{centred at $x$} has
$v(\mathfrak m)>0$. Unless otherwise stated, valuations are nontrivial
quasi-monomial real valuations. Discrepancy is extended linearly in the
weights on a log-smooth presentation. For every valuation $v$ over $x$,
put $v(0)=+\infty$ and, for $t\ge0$, set
\[
 \mathfrak a_t(v)=\{f\in R:v(f)\ge t\},\qquad
 \mathfrak a_{>t}(v)=\{f\in R:v(f)>t\}.
\]
The associated graded algebra and initial form are
\[
 \operatorname{gr}_v R=\bigoplus_{t\ge0}
       \mathfrak a_t(v)/\mathfrak a_{>t}(v),\qquad
 \operatorname{in}_v(f)=[f]_{v(f)}\quad(f\ne0).
\]
Degrees outside the value semigroup have zero summand. These conventions
also apply when the centre is a proper generization of $x$.
For a valuation centred at $x$, set
\[
 \vol(v)=\lim_{t\longrightarrow\infty}
       \frac{n!\,\operatorname{length}(R/\mathfrak a_t(v))}{t^n}.
\]
The limit exists by Cutkosky's theorem for graded families of
$\mathfrak m$-primary ideals \cite[Theorem~1.1 (=5.5)]{CutkoskyVolumes}.
Its hypothesis on the nilradical of the completion holds because $R$
is excellent and normal, hence has reduced completion.  The theorem
applies first to integral thresholds; squeezing between adjacent integral
thresholds gives the displayed real-parameter limit.
For a klt pair $(S,\Delta)$ and a centred valuation of finite log discrepancy,
the normalized volume is
\[
 \widehat{\vol}_{S,\Delta}(v)=A_{S,\Delta}(v)^n\vol(v).
\]
Its $n$th root is $A_{S,\Delta}(v)\vol(v)^{1/n}$, with the same minimizers.
Here and throughout the volume formulas, $n=\dim S$.

Fix a nonzero $h\in R$ such that $(S,H)$ is lc, where
$H=\divisor(h)$, and finitely many nonzero $F_j\in R$ and
$m_j\in\Z_{>0}$. The \emph{common lc slice} is
\begin{equation}\label{eq:E-slice}
 \mathcal Q=\{v:A(v)=v(h)=1,\quad v(F_j)\ge m_j\text{ for every }j\}.
\end{equation}
Assume that it contains a centred valuation. Its cone of rays is defined
by $A(v)=v(h)>0$ and $v(F_j)\ge m_jv(h)$.
An allowed height is
\begin{equation}\label{eq:E-height}
 b(v)=b_0v(h)-\sum_{\ell}b_{\ell}v(G_{\ell}),
 \qquad b_{\ell}\ge0,\quad G_{\ell}\in R\setminus\{0\},
 \qquad b|_{\mathcal Q}>0.
\end{equation}
On centred valuations for which $b(v)>0$, the degree-one homogeneous
height $b$ similarly gives the scale-invariant quantity $b(v)^n\vol(v)$ and its
$n$th root $b(v)\vol(v)^{1/n}$, with the same minimizers.
All constants in this section may depend on this individual germ and
its finite data. The constants in the strict volume gap below depend
only on dimension and the stated separation.

Our goal is to minimize this functional on the common lc rays and,
for rational heights, realize its minimizer as an ordinary
normalized-volume minimizer of an effective rational klt pair.
The realization must be exact: finite generation will be applied to
that pair and its actual minimizer. We first control retraction to
the lc locus and volume under interpolation; these estimates will
allow one rational perturbation to exclude every competing valuation
outside the common slice.

\subsection{Retractions, jets, and the topology of the slice}

Let $P\to S$ be a smooth model, $F$ a reduced snc divisor on $P$, and
$v$ a valuation with centre on $P$. Write $r_F(v)$ for its monomial
retraction using the components of $F$ through the centre, with their
$v$-weights. It is the trivial valuation if there are no such components.
Retraction decreases the value of every function regular on the chart
and preserves the values of its boundary monomials; see
\cite[Lemma~4.7, Corollary~4.8, and Proposition~5.1]{JonssonMustata}.
The following quantitative strengthening is needed here.

\begin{lemma}[Logarithmic jet estimate]\label{lem:E-log-jet}
Fix finitely many regular affine functions on $S$. On a fixed resolution
$P$, with $F$ reduced snc, there is a constant $C$ such that, for every
polynomial $g\ne0$ of degree at most $D$ in those functions and every
quasi-monomial $v$ over $x$,
\begin{equation}\label{eq:E-jet}
 r_F(v)(g)\le v(g)\le r_F(v)(g)+C(D+1)A_{P,F}(v).
\end{equation}
The same constant works when the centre is not closed and when some
weights disappear on passing to a face of the monomial cone.
\end{lemma}

\begin{proof}
Cover the relevant part of $P$ by finitely many smooth etale coordinate
charts. At the centre of $v$, write the coordinates of the components
of $F$ through that centre as $z_1,\ldots,z_s$ and the remaining
coordinates as $y_1,\ldots,y_{n-s}$. Thus $v(z_i)>0$.
Choose a closed specialization of the centre in this chart, avoiding
the other components of $F$, and translate the $y$-coordinates to
vanish there. Both $v$ and its retraction are nonnegative on this
closed-point local ring: the original centre generalizes this point.
A unit in that ring consequently has value zero for both valuations.
Put $a_i=v(z_i)$ and expand $g$ in the completed local ring.
The least $z$-weight is
\[
 d_0=\min\{\langle a,\alpha\rangle:
               [z^\alpha y^\beta]g\ne0\text{ for some }\beta\}.
\]
We first prove a degree bound independent of $a$.

The function field of the chart is a finite extension of
$\C(z,y)$. In a fixed basis, multiplication by each of the finitely
many affine functions has a matrix with fixed rational-function
entries. A polynomial of degree at most $D$ in these matrices has a
common denominator of degree $O(D+1)$ and numerator entries of the
same degree. Its characteristic equation, after clearing denominators,
gives
\begin{equation}\label{eq:E-algebraic-equation}
 \sum_{j=0}^{q}p_j(z,y)g^j=0,\qquad
 \deg p_j\le C_0(D+1),
\end{equation}
where $q$ is fixed and the coefficients are not all zero.
Translations of coordinates do not change this degree bound.
Give \emph{all} coordinates strictly positive rationally independent
weights. Let $\gamma$ be the leading exponent of $g$, and let $\delta_j$
be the leading exponent of a nonzero $p_j$. At least two distinct
powers, say $j$ and $k$, have the least weight in
\eqref{eq:E-algebraic-equation}. Independence of the weights gives
\[
 \delta_j+j\gamma=\delta_k+k\gamma.
\]
Consequently $|\gamma|\le C_1(D+1)$, after enlarging the constant.
This conclusion holds for every such choice of positive weights.

Now give the $z$-coordinates their fixed positive weights and give the
$y$-coordinates weight zero. There are finitely many $z$-exponents
of weight $d_0$, and only finitely many in a sufficiently small fixed
neighbourhood of this weight. Every vertex of the convex hull of the
tied $z$-exponents is exposed by a small generic perturbation of the
$z$-weights. Give the transverse coordinates sufficiently small positive
generic weights afterwards. The leading exponent projects to that
vertex: the positive gap to the other $z$-weights prevents an exponent
outside the least face from entering. The preceding bound therefore
bounds every vertex, and convexity bounds \emph{every} tied
$z$-exponent by $C_1(D+1)$. The same perturbation supplies at least one
pair $(\alpha,\beta)$ on the least face with
$|\alpha|+|\beta|\le C_1(D+1)$.

For this choice of $\alpha$, apply the product, over the $z$-coordinates,
of the operators
\[
 \prod_{\substack{0\le l\le C_1(D+1)\\ l\ne\alpha_i}}
             (z_i\partial_{z_i}-l).
\]
It kills every other tied $z$-exponent and multiplies the chosen
coefficient by a nonzero constant. Apply a transverse derivative
$\partial_y^\beta$ for a bounded exponent occurring in that coefficient.
Its new constant term is nonzero. These operators have total order
$O(D+1)$, and do not change any $z$-exponent. Their result $g'$ satisfies
\begin{equation}\label{eq:E-isolated-jet}
 g'=u z^\alpha\pmod{I_{>d_0}},\qquad u\text{ a local unit},
\end{equation}
where $I_{>d_0}$ is the monomial ideal generated by the $z$-monomials
of weight greater than $d_0$. This ideal is finitely generated.
The assertion descends from completion by faithful flatness. More
explicitly, faithful flatness first gives $g'=u z^\alpha+q$ with
$u$ regular and $q\in I_{>d_0}$. In the completion, $u$ differs from
the formal unit coefficient by an element of $(I_{>d_0}:z^\alpha)$.
This colon ideal is proper, because $z^\alpha$ has weight $d_0$,
and hence lies in the local maximal ideal. Thus $u$ is a unit.
Each monomial generator of $I_{>d_0}$ has value greater than $d_0$
for both $v$ and $r_F(v)$, and its regular coefficients have nonnegative
value. Thus
\[
 v(g')=r_F(v)(g')=d_0=r_F(v)(g).
\]
This also explains why a nonclosed centre causes no problem.

For completeness, a regular logarithmic vector field $L$ satisfies
\begin{equation}\label{eq:E-log-derivative}
 v(Lf)\ge v(f)-A_{P,F}(v).
\end{equation}
For a prime-order valuation, take a local logarithmic volume form
$\eta$ for $(P,F)$ and use
\[
 (Lf)\eta=f\,d\log(f)\wedge\iota_L\eta.
\]
At the prime on a higher smooth model, logarithmic coordinate forms
and $d\log(f)$ have nonnegative logarithmic order. Comparing the two
sides gives \eqref{eq:E-log-derivative}. Homogeneity gives it for rational
monomial weights, and rational approximation gives it for arbitrary
quasi-monomial weights. In this approximation the signs on a fixed
chart ring can be retained: the values of its finitely many generators
are rational piecewise linear functions on the monomial weight cone,
so one approximates in their common sign cells. Discrepancy and the
finitely many function values in question are continuous there.
The Euler operators and transverse derivatives used above are
logarithmic vector fields, and subtracting a scalar does not worsen
the estimate. Iterating \eqref{eq:E-log-derivative} and using
\eqref{eq:E-isolated-jet} proves the upper bound in \eqref{eq:E-jet}.
The lower bound is the retraction inequality.
\end{proof}

Suppose more generally that $\Psi$ is an effective lc rational sum of
Cartier divisors. Resolve its support and any further finite list of
ideals and divisors, and write
\[
 K_P+B_P=\pi^*(K_S+\Psi),\qquad F=B_P^{=1},\qquad
 e=A_{S,\Psi}(v),\qquad r=r_F(v).
\]
If $E$ is the reduced union of the support of $B_P$, the exceptional
locus, and the additional resolved divisors, then
\[
 e=A_{P,E}(v)+\sum_{D\subset E}(1-\operatorname{coeff}_D B_P)v(D).
\]
All summands are nonnegative. The positive numbers
$1-\operatorname{coeff}_D B_P$ for $D\not\subset F$ have a positive
minimum on this fixed resolution. Therefore
\begin{equation}\label{eq:E-retraction-errors}
 A_{P,F}(v)\le C e,\qquad
 |A_S(v)-A_S(r)|\le C e,
\end{equation}
and the evaluations of each resolved ideal or divisor differ at $v$
and $r$ by at most $Ce$. These statements follow also when $B_P$ has
negative coefficients: the displayed coefficients $1-\operatorname{coeff}_D B_P$
are still positive away from $F$. In particular, if $e=0$, then
Lemma~\ref{lem:E-log-jet} gives equality on every polynomial in the
fixed affine generators. Equality extends to the local ring and its
fraction field by localization. Hence $v=r$.

\begin{lemma}[Compactness and uniform centring]\label{lem:E-compactness}
The slice $\mathcal Q$ is a compact finite polyhedral complex of
monomial valuations on any log resolution adapted to $h$ and the
$F_j$. Evaluations of fixed functions are continuous on it.
For centred valuations with $A(v)=1$, there are constants $c,C>0$
depending on the germ such that
\begin{equation}\label{eq:E-centering}
 \vol(v)\ge c,\qquad
 \vol(v)\ge \frac{c}{v(\mathfrak m)}.
\end{equation}
Volume is continuous on the centred part of $\mathcal Q$.
\end{lemma}

\begin{proof}
Apply \eqref{eq:E-retraction-errors} to $\Psi=H$.
Every zero-discrepancy valuation is monomial on the coefficient-one
part of the fixed resolution. Conversely those monomial valuations
have zero discrepancy for $(S,H)$. Retain the strata whose closed
images contain $x$. This condition is preserved when a stratum is
replaced by a larger stratum, so the retained cones include their
faces. On a cone with primitive generators $E_i$ the equation
$A(v)=\sum_i a_iA(E_i)=1$ cuts out a compact simplex, because
$A(E_i)>0$. The additional inequalities are linear on an adapted
resolution. There are finitely many cones. Function evaluations are
monomial minima on them and agree on faces by localization to the
larger stratum, proving compactness and continuity.

We use Li's Izumi estimate \cite[Theorem~3.1]{LiIzumi}:
\begin{equation}\label{eq:E-izumi}
 v(f)\le C_I A(v)\ord_{\mathfrak m}(f)
 \quad(f\in R\setminus\{0\},\ v\text{ centred}).
\end{equation}
Thus, when $A(v)=1$, the valuation ideal $\mathfrak a_t(v)$ is contained
in $\mathfrak m^{\lceil t/C_I\rceil}$. Hilbert--Samuel growth gives the
first bound in \eqref{eq:E-centering}.
Choose fixed generators $z_1,\ldots,z_q$ of $\mathfrak m$, and let
$z=z_i$ have $v(z)=\delta=v(\mathfrak m)$. Since $R$ is a domain,
\[
 (\mathfrak a_t(v):z^j)=\mathfrak a_{t-j\delta}(v).
\]
Filtering $R/\mathfrak a_t(v)$ by the images of $z^jR$ gives
\[
 \operatorname{length}(R/\mathfrak a_t(v))
 =\sum_{0\le j<t/\delta}
   \operatorname{length}\bigl(R/(\mathfrak a_{t-j\delta}(v)+(z))\bigr).
\]
For $j\le t/(2\delta)$, Izumi bounds this summand below by
$\operatorname{length}(R/(\mathfrak m^{\lceil t/(2C_I)\rceil}+(z)))$.
The latter grows as a positive constant times $t^{n-1}$; this includes
$n=1$, when it is a positive constant. Since $z$ belongs to a fixed
finite list, the lower constant can be chosen uniformly. Divide by
$t^n/n!$ to obtain the second bound in \eqref{eq:E-centering}.

To prove continuity, consider valuations converging on one monomial
cone to a centred limit. Drop the components whose weights tend to
zero, and call the resulting retractions $r_i$. The retained positive
weights converge multiplicatively. On regular functions, the $r_i$
are therefore multiplicatively $1+o(1)$ equivalent to the limit.
Their maximal-ideal values, and those of the original valuations,
are bounded below. The reduced-log discrepancy measuring the dropped
weights tends to zero. Lemma~\ref{lem:E-log-jet} upgrades this fact to
a multiplicative comparison on \emph{all} regular functions as follows.

Write $q=\ord_{\mathfrak m}(f)\ge1$. With $c_0>0$ a common lower bound
for maximal-ideal values and $C_0q$ the common Izumi upper bound for
the values of $f$, fix $L$ with $Lc_0>2C_0$. Represent $f$ modulo
$\mathfrak m^{N}$, $N=\lceil Lq\rceil$, by a polynomial $p$ of degree
at most $N$ in fixed generators of $\mathfrak m$. Such a representative
exists also for localized regular functions by expanding their unit
denominators modulo $\mathfrak m^N$. The error has value greater than
both values of $f$, so $v_i(p)=v_i(f)$ and $r_i(p)=r_i(f)$.
The jet estimate bounds their difference by $o(1)q$, uniformly in $f$.
Since $r_i(f)\ge c_0q$, it follows that
$r_i\le v_i\le(1+o(1))r_i$ on all regular functions. Units have value
zero. Inclusion of valuation ideals now proves volume continuity.
Every convergent sequence has a subsequence in one of the finitely
many cones, which proves the assertion on the whole centred part.
\end{proof}

We record explicitly the uniform version of the truncation argument.
\begin{corollary}[Simultaneous approximation]\label{cor:E-all-functions}
In the situation of \eqref{eq:E-retraction-errors}, suppose
$A_S(v)$ and $A_S(r)$ are bounded above and
$v(\mathfrak m),r(\mathfrak m)\ge c_0>0$. Then a single constant $C$
works for every such $v$ and every $f\in R$:
\begin{equation}\label{eq:E-all-functions}
 r(f)\le v(f)\le(1+Ce)r(f).
\end{equation}
In particular,
$\vol(v)^{1/n}\ge(1+Ce)^{-1}\vol(r)^{1/n}$.
\end{corollary}
\begin{proof}
Use $N=\lceil L\ord_{\mathfrak m}(f)\rceil$ with the single constant
$L$ chosen in the preceding proof. Equations~\eqref{eq:E-jet} and
\eqref{eq:E-retraction-errors} give
$v(f)-r(f)\le C e\ord_{\mathfrak m}(f)\le Ce\,r(f)/c_0$.
Enlarge $C$. The ideal inclusions give the volume inequality.
\end{proof}

\subsection{Convex domination and strict volume comparison}

\begin{lemma}[Convex domination]\label{lem:E-convexity}
For $v_1,\ldots,v_k\in\mathcal Q$ and $p_s\ge0$ with
$\sum_s p_s=1$, there is $u\in\mathcal Q$ such that
\begin{equation}\label{eq:E-convexity}
 u(f)\ge\sum_s p_s v_s(f)\qquad(f\in R).
\end{equation}
If a positive-weight term is centred, then $u$ is centred. On the cone
of common lc rays, the same statement holds with
$A(u)=\sum_s p_sA(v_s)$, without requiring $A(v_s)=1$.
For an allowed height,
$b(u)\le\sum_s p_sb(v_s)$.
\end{lemma}

\begin{proof}
Fix finitely many functions $g_i$ whose averaged target values
$s_i=\sum_s p_sv_s(g_i)$ are positive, and choose positive rational
$s'_i\le s_i$. Include $h$ with the \emph{exact} target $1$ and each
$F_j$ with target $m_j$. Choose $N$ so that every $N/s'_i$ is integral,
and put $\mathfrak a=\sum_i(g_i^{N/s'_i})$. We claim that
$(S,\mathfrak a^{1/N})$ is not klt at $x$.
The multiplier-ideal summation formula on normal $\Q$-Gorenstein
complex varieties gives
\begin{equation}\label{eq:E-multiplier-sum}
 \mathcal J(\mathfrak a^{1/N})
 =\sum_{\substack{t_i\ge0\\\sum_i s'_it_i=1}}
       \mathcal J\Bigl(\prod_i(g_i)^{t_i}\Bigr).
\end{equation}
This is the mixed-ideal formula of
\cite[Theorem~3.2]{Takagi} and
\cite[Corollary~2 and Theorem~3.6]{JowMiller}; there is no smoothness
or Jacobian correction in this particular formula.
If the left side were the unit ideal at $x$, some summand would be the
unit ideal in the local ring. The resulting mixed pair would be klt,
so for every positive-weight $v_s$,
$\sum_i t_i v_s(g_i)<A(v_s)=1$. Averaging contradicts
\[
 \sum_i t_i\sum_s p_s v_s(g_i)=\sum_i t_is_i
       \ge\sum_i t_is'_i=1.
\]
The failure of klt is witnessed by a prime $E$ whose centre contains
$x$, and
$A(E)\le\min_i\{\ord_E(g_i)/s'_i\}$.
Thus $u=\ord_E/A(E)$ meets all the lower targets. In particular
$u(h)\ge1$, whereas lc of $(S,H)$ gives $u(h)\le A(u)=1$.
It is already in $\mathcal Q$.

Here is the geometric mechanism behind \eqref{eq:E-multiplier-sum}.
Resolve the summand ideals and their sum, with fixed divisors $D_i,D$,
and subdivide $\{p_i\ge0:\sum p_i=1/N\}$ by the rounding walls of
$\mathcal O_P(\lceil K_{P/S}-\sum_i p_iD_i\rceil)$.
The cellular complex indexed by chains of nonempty cells, with the
usual alternating face maps and inclusions on removing the first
cell, resolves
$\mathcal O_P(\lceil K_{P/S}-D/N\rceil)$.
In completed snc coordinates its exactness is coefficientwise:
the cells admitting a given Laurent monomial form a convex relatively
open subset of the exponent simplex, and their barycentric order
complex is contractible when nonempty. This subset is nonempty exactly
for monomials in the target. Indeed, locally one summand divided by
the invertible sum ideal is a unit, so its fixed divisor there equals
$D$. Local vanishing for each mixed multiplier-ideal term preserves
the surjective augmentation under pushforward. This gives the stated
sum formula in precisely the singular setting used above.

Let the rational lower targets increase to their averaged values.
Compactness from Lemma~\ref{lem:E-compactness} gives a point meeting
the finite list exactly as lower bounds. The closed sets imposing
these lower bounds, indexed by all regular functions, have the finite
intersection property. Compactness supplies a single $u$ satisfying
\eqref{eq:E-convexity} for every function. Positivity on generators of
$\mathfrak m$ proves the centring assertion. For unnormalized inputs,
put $a=\sum_s p_sA(v_s)$, apply the normalized assertion with weights
$p_sA(v_s)/a$ to $v_s/A(v_s)$, and multiply the result by $a$.
The formula for $b$ and the equality for $u(h)$ prove its height bound.
\end{proof}

\begin{corollary}[Opposite ratios]\label{cor:E-opposite-ratios}
Suppose rational functions $q,q'$ have regular numerators and
power-of-$h$ denominators, and $v(q')=-v(q)$ on $\mathcal Q$.
The interpolant in Lemma~\ref{lem:E-convexity} has
$u(q)=\sum_s p_sv_s(q)$ and $u(q')=\sum_s p_sv_s(q')$.
The assertion holds simultaneously for any number of such pairs.
\end{corollary}
\begin{proof}
Apply domination to both numerators and use the exact value of $h$.
The two lower inequalities have opposite sides and so must be equalities.
\end{proof}

\begin{lemma}[Volume interpolation and a strict gap]
\label{lem:E-volume-interpolation}
Let $w$ be centred, let $v$ be over $x$, and suppose
$u\ge tw+(1-t)v$ on $R$ for $0<t<1$. For $f\ne0$ with $w(f)>0$,
\begin{equation}\label{eq:E-volume-interpolation}
 \vol(u)\le t^{-n}
   \frac{tw(f)}{tw(f)+(1-t)v(f)}\vol(w).
\end{equation}
If $v$ is also centred, $t=1/2$, and
$v(f)\ge(1+\gamma)w(f)>0$ for $0<\gamma\le1$, then
\begin{equation}\label{eq:E-volume-gap}
 \vol(w)+\vol(v)-2\vol(u)\ge c(n,\gamma)\vol(w),
\end{equation}
where $c(n,\gamma)>0$ depends only on $n,\gamma$.
Consequently valuations whose volumes are within relative $o(1)$ of
the minimum on the same $A=1$ common slice are multiplicatively
$1+o(1)$ equal on \emph{all} regular functions. The conclusion is
uniform when the dimension is bounded. A sufficiently small fixed
relative error similarly gives any prescribed fixed multiplicative
closeness to one.
\end{lemma}

\begin{proof}
For a centred filtration, write its Laplace sum as the sum over its
jumps $a$ of
$\dim_\C(\mathfrak a_a/\mathfrak a_{>a})e^{-a/N}$.
Integration of the cumulative dimension function gives
\begin{equation}\label{eq:E-laplace}
 \lim_{N\longrightarrow\infty}N^{-n}
  \sum_a\dim_\C(\mathfrak a_a/\mathfrak a_{>a})e^{-a/N}
 =\vol(v).
\end{equation}
There are finitely many jumps below each bound, and polynomial growth
controls the tail.

Choose a homogeneous basis of
$\operatorname{gr}_wR/(\operatorname{in}_wf)$, lift it homogeneously,
and multiply the lifts by all nonnegative powers of
$\operatorname{in}_wf$. This is a basis of $\operatorname{gr}_wR$:
its independence follows from the domain property, and successive
division lowers the degree by $w(f)>0$, proving spanning.
Lift this basis to products $a_if^j$ in $R$. They span modulo arbitrary
high $w$-ideals, hence modulo high $u$-ideals since $u\ge tw$.
The $u$-value of $a_if^j$ is at least
$tw(a_i)+j(tw(f)+(1-t)v(f))$. Consequently its Laplace sum is at most
the $w$-Laplace sum at $N/t$, multiplied by
\[
 \frac{1-e^{-tw(f)/N}}
      {1-e^{-(tw(f)+(1-t)v(f))/N}}.
\]
Use \eqref{eq:E-laplace} to obtain
\eqref{eq:E-volume-interpolation}.

For the strict gap, fix $T>3$ and work in
$V_N=R/\mathfrak m^{\lceil CTN\rceil}$, where $C$ is large enough
for this fixed triple of centred valuations. The indicated power is
contained in all their threshold-$TN$ ideals. Cap both the $w$ and
$v$ flags at $TN$ and choose a basis simultaneously adapted to these
two flags; denote its capped values by $W_i,V_i$.
Lifts have $u$-value at least $(W_i+V_i)/2$. The identity
\[
 e^{-W_i/N}+e^{-V_i/N}-2e^{-(W_i+V_i)/(2N)}
     =\bigl(e^{-W_i/(2N)}-e^{-V_i/(2N)}\bigr)^2
\]
therefore bounds twice the capped $u$-Laplace sum above by the sum of
the other two, less the sum of these squares.

Choose a vector-space complement to
$\mathfrak a_{\gamma N/4}(w)$ in $R$ and multiply it by
$f^{\lfloor N/w(f)\rfloor}$. Every nonzero resulting element has
$w$-value less than $(1+\gamma/4)N$ and, for large $N$, has $v$-value
at least $(1+2\gamma/3)N$. Multiplication is injective, and this
subspace remains injective in $V_N$ by its upper $w$-value bound.
Its dimension is asymptotic to
\[
 \frac{\vol(w)}{n!}(\gamma N/4)^n.
\]
Its projection from the high-$v$ subspace modulo the high-$w$
subspace has this rank. Hence at least this many entries of a
simultaneously adapted basis satisfy
$W_i<(1+\gamma/4)N$ and $V_i\ge(1+2\gamma/3)N$.
Each contributes at least
\[
 \left(e^{-(1+\gamma/4)/2}-e^{-(1+2\gamma/3)/2}\right)^2
\]
to the preceding sum of squares. Divide by $N^n$, first let
$N\to\infty$, and then let $T\to\infty$.
The capped tail errors are bounded by polynomial functions of $T$
times $e^{-T}$ and disappear. This proves
\eqref{eq:E-volume-gap}, for example with any smaller positive
constant than
\[
 \frac{(\gamma/4)^n}{n!}
 \left(e^{-(1+\gamma/4)/2}-e^{-(1+2\gamma/3)/2}\right)^2.
\]

Finally interpolate two near-minimizers by
Lemma~\ref{lem:E-convexity}. Their interpolant has at least the
minimum volume. If their values on some regular function differed
by a fixed factor greater than one, apply \eqref{eq:E-volume-gap},
interchanging the valuations if necessary. It gives a fixed relative
loss, contrary to near-minimality. This proves the simultaneous
all-functions assertion and its fixed-error version.
\end{proof}

\subsection{Exact realization by a rational klt pair}

We have established compactness, a uniform comparison with monomial
retraction, and strict volume improvement when two valuations differ.
The remaining task is to encode the finite order constraints in an lc
boundary and then choose one rational perturbation for which every
ordinary minimizer lies exactly on its zero-discrepancy locus.  The
parameter must work for all competing valuations at once.

\begin{lemma}[A boundary encoding the common slice]\label{lem:E-slice-boundary}
Choose positive rational numbers $\theta_j$ with
$\sum_j\theta_j<1$. For sufficiently general scalars $a_j$, put
$E_j=\divisor(h^{m_j}+a_jF_j)$ and
\begin{equation}\label{eq:E-slice-boundary}
 \Psi=(1-\sum_j\theta_j)H+
                     \sum_j\frac{\theta_j}{m_j}E_j.
\end{equation}
Then $\Psi$ is effective and lc, and its nonzero zero-discrepancy
rays over the germ are exactly the common lc rays. On every such ray,
\begin{equation}\label{eq:E-pencil-order}
 v(h^{m_j}+a_jF_j)=m_jv(h).
\end{equation}
In particular, the coefficients $\theta_j$ may be chosen arbitrarily
small while retaining these conclusions.
\end{lemma}

\begin{proof}
Set $\mathfrak b_j=(h^{m_j},F_j)$. For any valuation,
\begin{align*}
 A(v)&-(1-\sum_j\theta_j)v(h)
               -\sum_j\frac{\theta_j}{m_j}v(\mathfrak b_j)\\
 &=A(v)-v(h)+\sum_j\theta_j
       \left(v(h)-\min\{v(h),v(F_j)/m_j\}\right)\ge0.
\end{align*}
Equality holds exactly on the common lc rays, together with the
trivial valuation. Resolve the ideals and choose the $E_j$ as
general members of their pencils. On the resolution the residual
free members meet the existing boundary with snc support and
contain none of its strata. Their coefficients $\theta_j/m_j$ are
less than one. They therefore introduce no new zero-discrepancy
strata, proving the assertion about $\Psi$.
On the coefficient-one monomial skeleton, each free residual member
has value zero because it contains no stratum. The value of its
pencil member is consequently the fixed-ideal value
$\min\{m_jv(h),v(F_j)\}=m_jv(h)$, proving
\eqref{eq:E-pencil-order}.
\end{proof}

\begin{theorem}[Exact common-slice minimizer]\label{thm:E-exact-minimizer}
For every allowed height $b$, the homogeneous functional
$b(v)\vol(v)^{1/n}$ on centred common lc rays has a unique minimizing
ray. Its $A=1$ representative depends continuously on the coefficients
of $b$ while the height stays strictly positive on the fixed slice.
It is fixed by every automorphism preserving the slice and the height.

If the coefficients of $b$ are rational, its minimizer $w$ is the
normalized-volume minimizer of an effective rational klt pair on
$x\in S$. More precisely, one can choose an effective rational lc
Cartier boundary $\Psi$ with exactly the common lc rays as its
zero-discrepancy locus, and one positive rational number $\lambda$,
fixed independently of the competing valuation, such that
\begin{equation}\label{eq:E-klt-perturbation}
 \Psi_\lambda=\Psi-\lambda b_0H
                       +\lambda\sum_\ell b_\ell\divisor(G_\ell)
\end{equation}
is effective and klt, and $w$ minimizes
$\widehat{\vol}_{S,\Psi_\lambda}$ among all centred valuations.
On the common lc rays its discrepancy is $\lambda b$.

For the same rational height, the algebra $\operatorname{gr}_wR$
is finitely generated and normal.
Its spectrum $S^0$ is a klt $\Q$-Gorenstein cone.
Write $h^0,F_j^0,G_\ell^0$ for initial forms, $w^0$ for the grading
valuation, and $\mathcal Q^0,b^0$ for the analogous central slice
and height. Then $(S^0,\divisor(h^0))$ is lc,
$A_{S^0}(w^0)=w^0(h^0)=1$, $w^0\in\mathcal Q^0$, and
$b^0$ is strictly positive on $\mathcal Q^0$.
The valuation $w^0$ minimizes $b^0\vol^{1/n}$ on its centred part.
\end{theorem}

\begin{proof}
\emph{Attainment, uniqueness, and continuity.}
Compactness gives $\min_{\mathcal Q}b>0$.
A minimizing sequence has bounded volume, so
\eqref{eq:E-centering} bounds its maximal-ideal values away from zero.
A subsequence converges in $\mathcal Q$ to a centred valuation.
Lemma~\ref{lem:E-compactness} gives volume continuity and hence
attainment. If two minimizing rays existed, normalize each by $b=1$.
Their midpoint has a dominating common lc interpolant $u$ with
$b(u)\le1$. If the minimum value of $b\vol^{1/n}$ is $M$, then
$\vol(u)\ge M^n/b(u)^n\ge M^n$.
The strict gap \eqref{eq:E-volume-gap} contradicts this unless the
two height-normalized valuations agree on every regular function.
Thus they are the same ray. For heights whose coefficients converge
to those of a positive height, the minimum heights stay bounded below
and one fixed centred competitor bounds the minimum values above.
The same compactness argument shows that every subsequential limit
of minimizers is the unique limiting minimizer. This proves continuity.
Automorphisms preserving the data preserve the functional, so
uniqueness proves equivariance.

\emph{Perturbing the boundary.}
Assume from now on that the height is rational. Choose $\Psi$ by
Lemma~\ref{lem:E-slice-boundary}, and resolve additionally all
$G_\ell$ and $\mathfrak m$.
For small rational $\lambda>0$, \eqref{eq:E-klt-perturbation} is
effective, since the coefficient of $H$ in \eqref{eq:E-slice-boundary}
is positive. It is klt near $x$: on each coefficient-one component
its new discrepancy is $\lambda b>0$, while the other finitely many
resolved discrepancies remain positive for sufficiently small
$\lambda$. The log-smooth criterion then tests every valuation.

\emph{Uniform control of every relevant competitor.}
Let $v$ be a centred quasi-monomial valuation with $A(v)=1$ and
\begin{equation}\label{eq:E-sublevel}
 A_{S,\Psi_\lambda}(v)\vol(v)^{1/n}\le\lambda M_0,
\end{equation}
where $M_0$ is the value of $b\vol^{1/n}$ at one fixed centred point
of $\mathcal Q$. This includes every normalized-volume minimizer of
the klt pair. Existence and quasi-monomiality for effective rational
klt pairs are supplied by \cite[Section~7]{Blum} and
\cite[Theorem~1.2]{XuQuasimonomial}.
The following constants are uniform over \emph{all} valuations in
\eqref{eq:E-sublevel} and all sufficiently small $\lambda$.

Put $e=A_{S,\Psi}(v)\ge0$ and retract $v$ to the coefficient-one
part of the fixed resolution, obtaining $r$.
Izumi bounds every fixed function value at $v$ when $A(v)=1$;
thus $|b(v)|\le C$. The first bound in
\eqref{eq:E-centering} and \eqref{eq:E-sublevel} give
\[
 e+\lambda b(v)\le C\lambda,\qquad e\le C\lambda.
\]
Equation~\eqref{eq:E-retraction-errors} yields
\begin{equation}\label{eq:E-uniform-retraction}
 A_S(r)=1+O(e),\qquad b(v)=b(r)+O(e).
\end{equation}
For small $\lambda$, $r$ is nontrivial and its normalized ray lies in
$\mathcal Q$. Strict positivity of $b$ on the compact slice therefore
gives uniform constants $\beta,B>0$ with
\[
 \beta\le b(r)\le B,\qquad b(v)\ge\beta/2.
\]
The sublevel inequality now bounds $\vol(v)$ above. The second bound
in \eqref{eq:E-centering} bounds $v(\mathfrak m)$ below; since
$\mathfrak m$ was resolved, its value at $r$ differs by $O(e)$.
Thus $r(\mathfrak m)$ is also bounded below and $r$ is an admissible
centred competitor. Corollary~\ref{cor:E-all-functions} supplies
\begin{equation}\label{eq:E-fixed-comparison}
 v(f)\le(1+Ce)r(f)\qquad\text{for every }f\in R,
\end{equation}
with a single $C$ independent of $f$, $v$, and small $\lambda$.

\emph{One fixed parameter forces exactness.}
Use \eqref{eq:E-uniform-retraction} and
\eqref{eq:E-fixed-comparison} to obtain
\begin{align}
 &(e+\lambda b(v))\vol(v)^{1/n}
       -\lambda b(r)\vol(r)^{1/n}\notag\\
 &\hspace{1cm}\ge
 \frac{e\,[1-\lambda C'-\lambda Cb(r)]}{1+Ce}
                 \vol(r)^{1/n}.\label{eq:E-exact-excess}
\end{align}
Choose \emph{one} positive rational $\lambda$ small enough for all
preceding bounds and for $\lambda(C'+CB)<1$.
This choice depends only on the fixed germ, boundary, and height;
it is independent of every competitor in \eqref{eq:E-sublevel} and
of every regular function. If such a competitor has $e>0$,
\eqref{eq:E-exact-excess} is strictly positive, so retraction strictly
improves its normalized-volume root functional.
An actual klt minimizer consequently has $e=0$, lies exactly on the
common lc slice, and is $w$ by uniqueness. Every competitor outside
\eqref{eq:E-sublevel} already has larger value than the fixed slice
competitor. This proves exact realization for this single rational
parameter; it is not merely a limiting assertion as $\lambda\to0$.
The argument applies for every smaller positive rational parameter
satisfying the same bounds.

\emph{Stable degeneration of this actual minimizer.}
We now apply stable degeneration to the effective rational klt pair
$(S,\Psi_\lambda)$ and its actual minimizer $w$.
By \cite[Theorems~1.1--1.2]{XuZhuangStable}, its associated graded is
finitely generated and defines a normal klt log Fano cone pair
$(S^0,\Psi_\lambda^0;w^0)$. The grading minimizes normalized volume
on the central pair; ordinary volume and pair discrepancy are
preserved, as in \cite[Lemma~4.10 and Theorem~3.5, arXiv version]{LiXuStable}.
The general uniqueness theorem for the klt-pair minimizer is
\cite[Theorem~1.1]{XuZhuangUnique}; the constrained uniqueness needed
here was proved before invoking it.

Each component used in \eqref{eq:E-klt-perturbation} is the divisor
of a regular function. For a principal ideal $(f)$ and a valuation
filtration, multiplicativity gives
$\operatorname{in}_w((f))=(\operatorname{in}_wf)$.
Equivalently its Rees homogenization, using its exact order, cuts a
Cartier divisor without a fibre component. Hence
$\Psi_\lambda^0$ is an effective rational sum of principal Cartier
divisors. Since $K_{S^0}+\Psi_\lambda^0$ is $\Q$-Cartier, so is
$K_{S^0}$, and removing the effective boundary shows that $S^0$ is klt.
Preservation of pair discrepancy and of all equation orders gives
\[
 A_{S^0}(w^0)=A_S(w)=1,\quad
 w^0(h^0)=1,\quad w^0(F_j^0)=w(F_j)\ge m_j.
\]

The choice of the $\theta_j$ can be made arbitrarily small. For each
such choice the exact minimizer is still the unique $w$, and thus
the graded algebra remains the same. Choose the corresponding
$\lambda$ smaller as necessary. The central effective klt boundaries
then contain $(1-o(1))\divisor(h^0)$.
For every prime over $S^0$, the resulting discrepancy inequalities
imply, on taking the coefficient limit,
$A_{S^0}(E)\ge\ord_E(h^0)$. This proves that
$(S^0,\divisor(h^0))$ is lc.

Finally take $v'\in\mathcal Q^0$. By \eqref{eq:E-pencil-order},
the initial equation of $h^{m_j}+a_jF_j$ is either $(h^0)^{m_j}$,
or $(h^0)^{m_j}+a_jF_j^0$ when $w(F_j)=m_j$.
In either case its $v'$-value is at least $m_jv'(h^0)$.
Therefore
\begin{equation}\label{eq:E-central-discrepancy}
 0<A_{S^0,\Psi_\lambda^0}(v')\le\lambda b^0(v'),
\end{equation}
with equality on the right at $w^0$.
The left inequality proves positivity of $b^0$ on the entire common
slice, including its noncentred faces. For centred $v'$, central
normalized-volume minimality and \eqref{eq:E-central-discrepancy}
give
\[
 \lambda b^0(v')\vol(v')^{1/n}
 \ge A_{S^0,\Psi_\lambda^0}(v')\vol(v')^{1/n}
 \ge\lambda b^0(w^0)\vol(w^0)^{1/n}.
\]
This is the required constrained minimality on the central slice.
\end{proof}

\begin{corollary}[Near-minimizers for a height]\label{cor:E-height-near-minimizers}
Fix an allowed height $b$ and normalize its common rays by $b=1$.
If $w_i,v_i$ have functional values within relative $o(1)$ of the
minimum, then $w_i=(1+o(1))v_i$ on all regular functions, with an error
independent of the function. The same conclusion, after $A=1$
normalization, holds for minimizers of heights whose coefficients
converge to those of $b$.
\end{corollary}
\begin{proof}
For the height-normalized pair, the midpoint interpolant has height
at most one. Its volume is therefore at least the minimum of volume
on the height-one slice. The proof using
\eqref{eq:E-volume-gap} applies verbatim: any fixed separation on one
function would force a fixed relative volume loss. For converging
heights $b_i$, positivity on the compact slice gives
$b_i/b=1+o(1)$ uniformly. Their minimizing rays, normalized by $b=1$,
are thus near-minimizers for $b$, and the preceding assertion applies.
Evaluation on $h$, which equals $A$ on the common rays, shows that
the factors needed to normalize by $A=1$ converge as well.
\end{proof}

\begin{corollary}[Domination by the minimizer]\label{cor:E-domination}
For the minimizing ray $w$ and every $v\in\mathcal Q$,
\begin{equation}\label{eq:E-domination}
 \frac{v(f)}{b(v)}\le 2^n\frac{w(f)}{b(w)}
                 \qquad(f\in R).
\end{equation}
For rational heights the same statement holds on the central slice
with all symbols carrying a superscript $0$.
\end{corollary}
\begin{proof}
Normalize both rays by $b=1$ and interpolate them with weights $1/2$.
The resulting centred $u$ has $b(u)\le1$, so minimality gives
$\vol(u)\ge\vol(w)$. For $w(f)>0$,
\eqref{eq:E-volume-interpolation} gives
\[
 1\le2^n\frac{w(f)}{w(f)+v(f)}.
\]
Thus $v(f)\le(2^n-1)w(f)$, which implies the stated bound.
If $f$ is a unit both values are zero. The central assertion follows
from the same argument and Theorem~\ref{thm:E-exact-minimizer}.
\end{proof}

\begin{lemma}[One-sided height tilt]\label{lem:E-one-sided-tilt}
Let $G\in R\setminus\{0\}$, $m\in\Z_{>0}$, and
$g(v)=v(G)/m-v(h)$. Suppose $g(v)/v(h)\ge2c$ on some common lc ray,
where $c>0$ is rational. Set
\[
 M=\max_{\mathcal Q}g\ge2c,\qquad
 \mu_*=(M-c)^{-1},\qquad
 b_\mu(v)=v(h)+\mu(cv(h)-g(v)).
\]
For $0\le\mu<\mu_*$ the height is positive. There is a rational
$\mu$ in this interval whose $A=1$ minimizing valuation $w$ satisfies
\[
 g(w)>0,\qquad b_\mu(w)\ge1.
\]
Moreover $b_\mu\le2+1/c$ on $\mathcal Q$, and the corresponding
central height obeys the same bound on $\mathcal Q^0$.
\end{lemma}

\begin{proof}
As $\mu\uparrow\mu_*$, the minimum of
$b_\mu\vol^{1/n}$ tends to zero. To prove this even when an endpoint
zero occurs at a noncentred valuation $v$, interpolate $v$ with one
fixed centred point $w_1\in\mathcal Q$, giving $w_1$ weight $t$.
For $0<\mu_*-\mu\le t$, the interpolant $u$ has
$b_\mu(u)=O(t)$. Since $v(h)=w_1(h)=1$, the improved volume estimate
\eqref{eq:E-volume-interpolation}, applied with $f=h$, gives
$\vol(u)\le t^{1-n}\vol(w_1)$.
Hence $b_\mu(u)\vol(u)^{1/n}=O(t^{1/n})$.
The lower volume bound in \eqref{eq:E-centering} then forces the
minimizing height to tend to zero. In particular its minimizing
valuation eventually has $g(w)>c$.

If the minimizer at $\mu=0$ has $g(w)>0$, take $\mu=0$.
Otherwise continuous dependence in
Theorem~\ref{thm:E-exact-minimizer} gives a nonempty open interval
of parameters for which $0<g(w)<c$. Choose a rational parameter
in that interval. Then $b_\mu(w)=1+\mu(c-g(w))\ge1$.
Finally $v(G)\ge0$ gives $g(v)\ge-1$ on the normalized slice, while
$\mu<\mu_*\le1/c$. Thus
$b_\mu(v)\le1+(c+1)/c=2+1/c$.
The same calculation applies centrally because $G^0$ is regular and
$v'(h^0)=1$ on $\mathcal Q^0$.
\end{proof}
